% Litz project adaptation of Michael Shell's bare_jrnl.tex, IEEEtran 1.8b.
% Upstream original and LPPL notice are preserved in references/templates.
% This short specimen checks typesetting only; no physical solver is run.
\documentclass[journal,10pt]{IEEEtran}
\usepackage{amsmath,amssymb}
\usepackage{newtxmath}
\usepackage[no-math]{fontspec}
\setmainfont{Times New Roman}
\usepackage{xeCJK}
\setCJKmainfont[AutoFakeBold=2]{SimSun}
\usepackage{booktabs}
\usepackage{graphicx}
\usepackage[hidelinks]{hyperref}
% Keep both independent panels at 88.9 mm while retaining the class margins.
\setlength{\columnsep}{\dimexpr\textwidth-177.8mm\relax}
\title{Hierarchical Verification and Common-Terminal Current Sharing in Litz-Wire Models}
\author{Litz Wire Modeling Project}
\begin{document}
\maketitle
\begin{abstract}
This specimen exercises the IEEE journal layout using frozen project results.
It distinguishes analytical calibration, common-terminal current allocation,
and finite-length longitudinal PEEC. The purpose is to verify the compilation
chain and typography before further manuscript revision.
\end{abstract}
\begin{IEEEkeywords}
Litz wire, common-terminal constraint, PEEC, numerical verification.
\end{IEEEkeywords}
\section{Scope and Conventions}
All phasors are RMS quantities with the convention $\exp(\mathrm{j}\omega t)$.
The baseline uses $f=200\,\mathrm{kHz}$, $I_{\mathrm{t}}=20\,\mathrm{A}$,
and $\sigma=5.8\times10^7\,\mathrm{S/m}$ at $20\,{}^{\circ}\mathrm{C}$.
The $2\,\mathrm{mm}$ round wire is a calibration specimen; its DC current
density exceeds the original design limit and it is not a compliant final bundle.

\textbf{中文字体短测：}沿途绝缘、端部并联；采用共同端电压与总复电流约束。
几何展示、解析预测、二维场解和实际三维验证分别标明。
宋体用于中文，Times New Roman 用于西文与数字。

\section{Common-Terminal Model}
The strand currents are unknowns rather than prescribed equal currents:
\begin{equation}
\begin{bmatrix}\mathbf Z&-\mathbf1\\\mathbf1^{\mathsf T}&0\end{bmatrix}
\begin{bmatrix}\mathbf i\\V\end{bmatrix}
=\begin{bmatrix}-\mathbf e\\I_{\mathrm t}\end{bmatrix}.
\label{eq:common}
\end{equation}
Within the specified linear model, equal sharing requires
\begin{equation}
\mathbf Q\left[\frac{I_{\mathrm t}}{N}\mathbf Z\mathbf1+\mathbf e\right]
=\mathbf0,\quad
\mathbf Q=\mathbf I_N-\frac{\mathbf1\mathbf1^{\mathsf T}}N.
\label{eq:sharing}
\end{equation}
Equation~\eqref{eq:sharing} is a model-scoped algebraic criterion and does not
establish a ranking of real transposed geometries. DTM-based strand permutation
provides a relevant comparison framework~\cite{cruciani2023}.

\section{Frozen Numerical Check}
Table~\ref{tab:round} uses the saved Bessel calibration values. RMS normalization
gives $P=|I|^2R_{\mathrm{ac}}$ without a factor of $1/2$.
Figure~\ref{fig:round} reuses the independently exported Wolfram PDF without
replotting. Its physical width is $88.9\,\mathrm{mm}$.
\begin{figure}[!t]
\centering
\includegraphics[width=88.9mm]{../outputs/publication_20261006_full_05/panels/en/q1_magnitude.pdf}
\caption{Frozen Bessel radial current-density magnitude for the calibration
wire. The field convention and RMS normalization are unchanged.}
\label{fig:round}
\end{figure}
\begin{table}[!t]
\caption{Round-wire analytical calibration at 200 kHz}
\label{tab:round}
\centering
\begin{tabular}{lr}
\toprule
Quantity & Frozen value\\
\midrule
Skin depth (mm)&0.147771655\\
$R_{\mathrm{dc}}$ ($\mathrm{m\Omega/m}$)&5.488101486\\
$R_{\mathrm{ac}}$ ($\mathrm{m\Omega/m}$)&20.016203628\\
$R_{\mathrm{ac}}/R_{\mathrm{dc}}$&3.647199980\\
Loss ($\mathrm{W/m}$)&8.006481451\\
\bottomrule
\end{tabular}
\end{table}

\clearpage
\section{Discussion and Conclusions}
The longitudinal PEEC implementation assumes a fixed axial current basis;
transverse loops and true end-region redistribution remain outside its scope.
RT0 imposed-field loss is distinct from transport resistance. Preserved
A--$\phi$ field samples retain their failed potential and gauge acceptance
status. A successful typesetting test does not change those physical outcomes.
Figure~\ref{fig:geometry} shows archived geometry alone, with no claimed
electromagnetic performance ranking.
\begin{figure}[!t]
\centering
\includegraphics[width=88.9mm]{../outputs/publication_20261006_full_05/panels/en/multilevel_detail.pdf}
\caption{An actual 3 mm detail of the archived 24-strand multilevel
geometry. Copper diameter: 0.100 mm. Highlighted physical strand identities
are retained.}
\label{fig:geometry}
\end{figure}

\bibliographystyle{IEEEtran}
\bibliography{references}
\end{document}
